\documentclass[10pt,a4paper]{amsart}
\usepackage[margin=19mm]{geometry}
\usepackage{amsmath,amssymb,mathtools}
\usepackage[scr=rsfs,cal=euler]{mathalfa}
\usepackage{xfrac}
\usepackage[unicode,hidelinks,bookmarks=false]{hyperref}
\newcommand{\ie}{i.e.\ }
\newcommand{\cf}{cf.\ }
\newcommand{\resp}{respectively\ }
\newcommand{\Subsection}{\subsection}
\providecommand{\nobreakdash}{\nobreak\hbox{-}}
\setlength{\emergencystretch}{2em}
\multlinegap0pt
\usepackage{amssymb}
\usepackage{mathtools}
\usepackage{tikz}
\usepackage{xcolor}
\usepackage[shortlabels]{enumitem}
\newtheorem{lemma}{Lemma}
\newcommand{\abs}[1]{\left\lvert#1\right\rvert}
\newcommand{\symdiff}{\mathbin{\triangle}}    % matching symmetric difference (source: \mathbin{\triangle})
\title{Matching complements in subcubic graphs\\[2pt] and a proof of the 3-Decomposition Conjecture}
\author{Jicheng Ma}
\date{Source: arXiv:2608.25385v2 (CC BY 4.0)}
\begin{document}
\maketitle

\noindent\emph{Source and licence.} Matching complements in subcubic graphs\\[2pt] and a proof of the 3-Decomposition Conjecture. Version arXiv:2608.25385v2 (CC BY 4.0), distributed by arXiv under the Creative Commons Attribution 4.0 International licence, http://creativecommons.org/licenses/by/4.0/, as recorded in the arXiv licence metadata for this exact version and shown on the abstract page https://arxiv.org/abs/2608.25385v2. Full licence notice: \texttt{licenses/arxiv-2608.25385-CC-BY-4.0.txt}.\par\smallskip

\noindent\textit{Editorial scope.} Counted unit: the article's lemma lem:first-entry (source0.tex lines 673--683). The article's complete original proof of it is delivered with it (source0.tex lines 685--704, one \texttt{proof} environment of 20 lines). No mathematics is written, repaired or re-derived: the statement and the proof are the article's own lines, in the article's own notation, and no retained line is substituted or re-flowed.  No further source-local context is needed: every cross-reference of every retained line resolves inside the two delivered spans.  Nothing was omitted from any retained span, so the package carries no omission record.  No retained line contains a citation, so the wrapper carries no bibliography and no bibliography entry.  No retained line contains a figure environment, an image-inclusion command or drawing code, so no image is delivered and the package declares no asset dependency.  Editorial formatting only, in addition: an AMS \texttt{amsart} preamble that restates the article's own declarations and macros used by the retained lines, namely the environments \texttt{lemma} on separate counters, together with the standard packages those lines need.  The retained environments print in this document as Lemma 1 in order of appearance, whereas the article prints the same items with the numbering of its own full document.  The complete native source of the article as distributed in the arXiv e-print for this exact version is bundled unchanged as \texttt{sources/208575/source0.tex}.
\begin{lemma}\label{lem:first-entry}
  Let $M,N$ be matchings of a finite loopless graph $G$, let $Q$ be a component
  of $G-M$, and let $x\notin V(Q)$. Assume that
  \begin{itemize}
    \item $x$ is exposed by $M$ and covered by $N$;
    \item every other vertex exposed by $M$ belongs to $Q$; and
    \item every vertex exposed by $N$ belongs to $Q$.
  \end{itemize}
  Then there is a matching $M'$ with $\abs{M'}=\abs{M}$ such that the component
  of $G-M'$ containing $Q$ contains every vertex exposed by $M'$.
\end{lemma}

\begin{proof}
  In $M\symdiff N$, the vertex $x$ has degree one, and its incident edge belongs
  to $N\setminus M$. Its component is therefore a simple alternating path. The
  other endpoint is exposed by exactly one of $M,N$, and hence belongs to $Q$ by
  the hypotheses.

  Traverse the path from $x$, and stop at its first vertex $d\in V(Q)$. The edge
  that enters $Q$ must belong to $M\setminus N$: if it lay outside $M$, it would
  be an edge of $G-M$ joining its preceding outside vertex to the component $Q$,
  contrary to the choice of $Q$. Thus the prefix from $x$ to $d$ begins with an
  $N\setminus M$ edge and ends with an $M\setminus N$ edge, and it contains
  equally many edges of the two matchings.

  Flip membership along this prefix. The result $M'$ is a matching of the same
  size as $M$; it covers $x$ and exposes $d$. The prefix meets $Q$ for the first
  time at $d$, so no complement edge internal to $Q$ is removed, and the final
  $M$-edge becomes a complement edge entering $Q$. Every old $M$-exposure except
  $x$ remains in $Q$, and the only new exposure is $d\in Q$. Hence all
  $M'$-exposures lie in the component containing $Q$.
\end{proof}

\end{document}
